\documentclass[11pt,a4paper,oneside,openright]{book}
\usepackage[spanish,activeacute,es-noshorthands]{babel}
\usepackage[utf8]{inputenc}
\usepackage[T1]{fontenc}
\usepackage{lmodern}
\usepackage{amsmath, amssymb, amsthm}
\usepackage{graphicx}
\usepackage[dvipsnames]{xcolor}
\usepackage{listings}
\usepackage{geometry}
\geometry{margin=2.5cm, headheight=14pt}
\usepackage{microtype}
\usepackage{booktabs}
\usepackage{enumitem}
\usepackage{fancyhdr}
\usepackage{tcolorbox}
\tcbuselibrary{skins, breakable}
\usepackage{caption}
\captionsetup{font=small, labelfont=bf}
\usepackage{hyperref}
\hypersetup{
  colorlinks=true, linkcolor=NavyBlue, citecolor=ForestGreen, urlcolor=BrickRed,
  pdftitle={M10 --- Compute Module 5 + carrier para INTISAT},
  pdfauthor={INTISAT}
}

\newtheorem{definicion}{Definicion}[chapter]
\newtheorem*{nota}{Nota}

\definecolor{codebg}{RGB}{248, 249, 250}

\lstdefinestyle{shell}{
  basicstyle=\ttfamily\footnotesize,
  backgroundcolor=\color{codebg},
  frame=single, rulecolor=\color{gray!30},
  framesep=4pt, breaklines=true,
  showstringspaces=false,
  literate={a}{{a}}1 {e}{{e}}1 {i}{{i}}1 {o}{{o}}1 {u}{{u}}1 {n}{{n}}1
}
\lstset{style=shell}

\newtcolorbox{intuibox}[1][]{
  colback=green!5, colframe=green!40!black,
  fonttitle=\bfseries, title=Intuicion,
  breakable, sharp corners=south, #1
}
\newtcolorbox{calculo}[1][]{
  colback=gray!4, colframe=gray!55,
  fonttitle=\bfseries, title=Calculo,
  breakable, sharp corners=south, #1
}
\newtcolorbox{payload}[1][]{
  colback=blue!4, colframe=blue!50!black,
  fonttitle=\bfseries, title=Aplicacion al INTISAT,
  breakable, sharp corners=south, #1
}
\newtcolorbox{cuidado}[1][]{
  colback=red!5, colframe=red!50!black,
  fonttitle=\bfseries, title=Cuidado,
  breakable, sharp corners=south, #1
}
\newtcolorbox{ejercicio}[1][]{
  colback=orange!5, colframe=orange!60!black,
  fonttitle=\bfseries, title=Ejercicios,
  breakable, sharp corners=south, #1
}

\pagestyle{fancy}
\fancyhf{}
\fancyhead[LE,RO]{\thepage}
\fancyhead[RE]{\nouppercase{\leftmark}}
\fancyhead[LO]{\nouppercase{\rightmark}}

\title{\Huge\bfseries M10 --- Compute Module 5\\[0.4em]
\Large Diseño del carrier board para el INTISAT\\[0.6em]
\large Manual de la coleccion CubeSat (M10 / 11)}
\author{Coleccion CubeSat para el proyecto INTISAT}
\date{\today}

\begin{document}

\frontmatter
\maketitle

\chapter*{Prefacio}
\addcontentsline{toc}{chapter}{Prefacio}

Este manual es la guia practica para migrar tu payload de Raspberry Pi 5
de desarrollo al \textbf{Compute Module 5} (CM5) montado en un carrier
board propio diseñado para el INTISAT.

El CM5 es la version "industrial" del Pi 5: mismo SoC BCM2712, mismo
software, pero compacto, sin USB / Ethernet expuestos, con eMMC en lugar
de SD, y con un conector mezzanine de 200 pines que vos llevas al
carrier que diseñes.

\section*{Audiencia}

\begin{itemize}[leftmargin=2em]
  \item Ingenieria electronica del payload.
  \item Diseñadores de PCB (KiCad, Altium, EAGLE).
  \item Integradores mecanicos del CubeSat.
\end{itemize}

\section*{Prerequisitos}

\begin{itemize}[leftmargin=2em]
  \item Manual M0 leido (conceptos de CubeSat).
  \item Manual de Electronica Practica (Capitulos 1-12).
  \item Manual Raspberry\_Pi\_5\_Profundo (capitulos 1-7).
  \item Familiaridad con KiCad o equivalente.
\end{itemize}

\section*{Objetivo del manual}

Que al final tengas:
\begin{enumerate}[leftmargin=2em]
  \item Diseño esquematico del carrier (KiCad).
  \item Layout del PCB de 4 layers.
  \item BOM completa con part numbers.
  \item Procedimiento de bring-up.
  \item Plan de test ambiental.
\end{enumerate}

\tableofcontents

\mainmatter

% =====================================================================
\chapter{Por que migrar a CM5}

El Raspberry Pi 5 retail es excelente para \textbf{desarrollo} en lab,
pero tiene problemas para vuelo:

\begin{itemize}[leftmargin=2em]
  \item \textbf{Tamaño}: 85$\times$56 mm. Casi todo un slot de 1U.
  \item \textbf{Conectores expuestos}: USB-C, USB-A x4, HDMI, Ethernet,
        GPIO header. Cada conector es un riesgo mecanico y de
        contaminacion.
  \item \textbf{SD card}: el slot mecanico de la SD se afloja con
        vibracion. La SD misma se corrompe con radiacion y temperatura.
  \item \textbf{Componentes no usados consumen}: WiFi, Bluetooth,
        Ethernet PHY... 0.5-1 W de consumo innecesario.
\end{itemize}

\section{Lo que aporta el CM5}

\begin{center}
\begin{tabular}{p{4cm}p{5cm}p{5.5cm}}
\toprule
\textbf{Aspecto} & \textbf{Pi 5 retail} & \textbf{CM5} \\
\midrule
Dimensiones        & 85$\times$56 mm    & 55$\times$40 mm (45\% mas chico) \\
Conector           & headers + USB + Ethernet & mezzanine DF40 200-pin \\
Almacenamiento     & SD card            & eMMC integrada (4-64 GB) \\
USB                & 4 puertos          & solo en pines, no conectores fisicos \\
Ethernet           & RJ-45              & solo PHY (vos pones el conector) \\
HDMI               & si                 & expuesto via mezzanine \\
Conector camara    & FFC 22-pin         & via mezzanine \\
Consumo            & 3-8 W              & 2.5-7 W (10\% menos) \\
Long-term avail.   & 5-7 años           & 10+ años (industrial grade) \\
Precio             & \$80                & \$50 + carrier custom \\
\bottomrule
\end{tabular}
\end{center}

\section{Lo que NO aporta el CM5}

\begin{itemize}[leftmargin=2em]
  \item NO es rad-hard. Sigue siendo COTS comercial.
  \item NO esta certificado para temperatura industrial extendida.
  \item NO tiene ECC en RAM.
  \item NO viene con carrier; tenes que diseñarlo.
\end{itemize}

\begin{payload}
Para el INTISAT (mision academica LEO de 1-2 años), el CM5 es \textbf{el
sweet spot}: tiene los beneficios mecanicos y de eMMC del industrial,
pero al costo del comercial. La radiacion en LEO baja no llega a niveles
que requieran rad-hard.
\end{payload}

% =====================================================================
\chapter{Anatomia del CM5}

\section{Especificaciones}

\begin{center}
\begin{tabular}{ll}
\toprule
\textbf{Spec} & \textbf{Valor} \\
\midrule
SoC                 & Broadcom BCM2712 \\
CPU                 & 4x ARM Cortex-A76 @ 2.4 GHz \\
GPU                 & VideoCore VII \\
RAM (variantes)     & 2 / 4 / 8 / 16 GB LPDDR4X \\
eMMC (variantes)    & Lite (0) / 4 / 16 / 32 / 64 GB \\
WiFi/BT (opcional)  & Cypress CYW43455 dual-band \\
Voltaje principal   & 5 V \\
Voltaje IO          & 3.3 V (configurable a 1.8 V) \\
Dimensiones         & 55 $\times$ 40 $\times$ 4.7 mm \\
Conector            & 2x Hirose DF40C-100DS (100 pin c/u $=$ 200 pin total) \\
Stack height        & 1.5 mm (al carrier) \\
Temperatura         & 0 a +85 °C (comercial) \\
\bottomrule
\end{tabular}
\end{center}

\section{Variantes para elegir}

\begin{center}
\begin{tabular}{lllll}
\toprule
\textbf{Variante} & \textbf{RAM} & \textbf{eMMC} & \textbf{WiFi} & \textbf{Precio} \\
\midrule
CM5 Lite (2 GB)    & 2 GB & --- & no & 25 USD \\
CM5 (4 GB / 32GB)  & 4 GB & 32 GB & no & 60 USD \\
CM5 (8 GB / 64 GB) & 8 GB & 64 GB & no & 95 USD \\
CM5 + WiFi (8/64)  & 8 GB & 64 GB & si & 105 USD \\
\bottomrule
\end{tabular}
\end{center}

\begin{payload}
\textbf{Recomendacion para INTISAT}: CM5 \textbf{4 GB / 32 GB sin WiFi}.
Tu pipeline IA cabe en 4 GB. eMMC de 32 GB alcanza para sistema operativo,
codigo, y buffer de scans (8-10 dias antes de overwrite). WiFi se desactiva
para vuelo, no vale la pena pagarlo.
\end{payload}

\section{Los 200 pines del mezzanine}

Los 200 pines del CM5 se dividen en:

\begin{itemize}[leftmargin=2em]
  \item \textbf{Power}: 14 pines (5V in, 3V3, GND).
  \item \textbf{GPIO}: 40 (mismos que el header de la Pi 5).
  \item \textbf{Buses dedicados}: 2x I2C, 5x SPI, 6x UART, USB 3.0 x4, USB 2.0 x4.
  \item \textbf{CSI-2 cameras}: 4 lanes diferenciales x2 instancias.
  \item \textbf{DSI displays}: idem.
  \item \textbf{HDMI}: 1 puerto.
  \item \textbf{Ethernet}: pines del PHY.
  \item \textbf{PCIe Gen 2 x1}.
  \item \textbf{Control}: nRUN (reset), nPOWER\_BTN, EEPROM\_nDIS, LED.
\end{itemize}

El pinout completo esta en el datasheet:
\url{https://datasheets.raspberrypi.com/cm5/cm5-datasheet.pdf}.

% =====================================================================
\chapter{Que debe tener el carrier}

\section{Lista de bloques funcionales}

El carrier para el INTISAT necesita:

\begin{enumerate}[leftmargin=2em]
  \item Conectores mezzanine DF40 (2 unidades hembra de 100 pin).
  \item Regulacion de 5V con proteccion (LDO o switching).
  \item Capacitores de decoupling (100 nF cerca de cada chip).
  \item Diodo Schottky de proteccion polaridad inversa.
  \item Fusible PTC reseteable.
  \item Level shifter TXS0108E para I2C externo.
  \item Sensor INA219 de corriente.
  \item Conector CSI-2 para OV5647.
  \item Conector SPI / GPIO para OLED SSD1351.
  \item Conector PC-104 para integracion al INTISAT.
  \item Boton de reset y LED de status.
  \item Pads de prueba para todos los buses.
  \item Conector de programacion de eMMC (rpiboot).
\end{enumerate}

\section{Diagrama de bloques}

\begin{center}
\begin{tabular}{|p{14cm}|}
\hline
\textbf{Bloque} \\
\hline
\textbf{Entrada de potencia}: PC-104 5V\_PAYLOAD $\to$ fusible PTC 5A $\to$
diodo Schottky $\to$ LDO 5V/6A $\to$ 5V interno. \\
\hline
\textbf{Mezzanine CM5}: 2x DF40C-100DS, stack height 1.5 mm. Conectados a
todas las lineas relevantes. \\
\hline
\textbf{Comms con OBC}: pines I2C del CM5 $\to$ TXS0108E $\to$ pines del PC-104. \\
\hline
\textbf{Monitor de potencia}: INA219 entre 5V externo y CM5, conectado a
I2C secundario del CM5. \\
\hline
\textbf{Camara}: conector FFC 15-pin a OV5647 directo del CSI-2 del CM5. \\
\hline
\textbf{OLED}: conector header 5-pin (MOSI, SCLK, CE, DC, RST + 3V3 + GND)
a SSD1351. \\
\hline
\textbf{Status}: LED bicolor controlado por GPIO. \\
\hline
\end{tabular}
\end{center}

% =====================================================================
\chapter{Diseño esquematico}

\section{Herramienta: KiCad}

KiCad es open-source, multi-plataforma, y suficiente para diseñar un PCB
de 4 layers de calidad aeroespacial. Tambien valen Altium (caro) o EAGLE
(deprecated).

\section{Hierarchical schematic}

El esquematico se divide en \textbf{hojas jerarquicas}:

\begin{enumerate}[leftmargin=2em]
  \item \texttt{TOP.kicad\_sch} --- vista general con bloques.
  \item \texttt{power.kicad\_sch} --- entrada PC-104, regulacion, INA219.
  \item \texttt{cm5.kicad\_sch} --- conectores DF40, decoupling.
  \item \texttt{io\_obc.kicad\_sch} --- TXS0108E + PC-104.
  \item \texttt{io\_camera.kicad\_sch} --- conector CSI-2.
  \item \texttt{io\_oled.kicad\_sch} --- conector SPI + GPIO.
  \item \texttt{debug.kicad\_sch} --- LED, boton, test pads.
\end{enumerate}

\section{Reglas de diseño electricas}

\begin{itemize}[leftmargin=2em]
  \item Todo pin de power tiene un capacitor 100 nF a menos de 3 mm.
  \item Capacitor de bulk 10 uF cerca de cada rail.
  \item Cada chip tiene su propio decoupling.
  \item Ground plane continuo en la layer 2.
  \item Power plane en layer 3 (dividido por rails).
  \item Signal routing en layers 1 y 4.
\end{itemize}

% =====================================================================
\chapter{Regulacion de potencia}

\section{De PC-104 a CM5}

El CM5 acepta 5V $\pm$ 5\%. La entrada del PC-104 puede tener picos.
Tres opciones:

\subsection{Opcion A: LDO}

\begin{itemize}[leftmargin=2em]
  \item \textbf{Pros}: simple, sin ruido switching, costo bajo.
  \item \textbf{Contras}: ineficiente; disipa $(V_{in} - V_{out}) \times I$ como calor.
\end{itemize}

Si $V_{in} = 5$ V y $V_{out} = 5$ V (con dropout de $\sim 0.2$ V), no hay
mucha disipacion. Pero el LDO solo "limpia" el rail, no convierte.

\subsection{Opcion B: Buck switching}

\begin{itemize}[leftmargin=2em]
  \item \textbf{Pros}: eficiencia $> 92\%$.
  \item \textbf{Contras}: ruidoso (switching freq 500 kHz - 2 MHz), mas
        complejo de diseñar.
\end{itemize}

\subsection{Opcion C: $\mu$Module}

Chips como el LTM4631 de Analog Devices vienen completos (regulador +
inductor + capacitores en un solo chip). Mas caros pero \emph{plug and play}.

\begin{payload}
Recomendacion: \textbf{$\mu$Module} (LTM4631 o similar). Ocupa $9\times 9$ mm,
da 5V $@$ 6A con $94\%$ eficiencia, y diseño simple. Cuesta USD 25 pero
ahorra horas de diseño y debugging.
\end{payload}

\section{Proteccion polaridad inversa}

Diodo Schottky (1N5817 o similar) en serie con la entrada:
\begin{itemize}[leftmargin=2em]
  \item Cae 0.3 V (Schottky tiene menos drop que diodo Si normal).
  \item Si conectan al reves, bloquea.
\end{itemize}

\section{Fusible PTC}

Fusible reseteable de 5A. Si hay un corto:
\begin{itemize}[leftmargin=2em]
  \item PTC sube su resistencia rapidamente.
  \item Corta el rail (limita la corriente).
  \item Cuando se enfria, vuelve a conducir.
\end{itemize}

% =====================================================================
\chapter{Subsistemas auxiliares}

\section{INA219 --- monitor de corriente}

Como vimos en el manual de electronica practica, el INA219 mide $V$ e $I$
del rail 5V del CM5. Lo conectas:

\begin{itemize}[leftmargin=2em]
  \item \textbf{IN+}: lado de la fuente.
  \item \textbf{IN-}: lado del CM5.
  \item \textbf{SDA, SCL}: I2C secundario del CM5 (no el que va al OBC).
  \item \textbf{V+}: alimentacion 3.3V del CM5.
\end{itemize}

Address default: 0x40. Configurable con jumpers a 0x41, 0x44, 0x45.

\section{TXS0108E --- level shifter}

Si el OBC del INTISAT usa 5V en su I2C y el CM5 usa 3.3V, necesitas un
level shifter bidireccional. El TXS0108E es la opcion estandar:

\begin{itemize}[leftmargin=2em]
  \item 8 canales (usamos 2: SDA y SCL).
  \item Auto-direction (no necesita pin de control).
  \item Voltajes: 1.65V a 5.5V en cada lado.
  \item Encapsulado: TSSOP-20.
\end{itemize}

\section{Conector CSI-2 para OV5647}

El CM5 expone 2 conectores de camara CSI-2 de 4 lanes cada uno. La OV5647
usa 2 lanes. Conectar via FFC 15-pin estandar del Pi.

\begin{nota}
Lamentablemente el CM5 expone CSI-2 via las lineas del mezzanine, no
directo. Hay que rutear las lineas diferenciales con $90~\Omega$ de impedancia
controlada en el carrier.
\end{nota}

\section{Conector OLED}

Header simple de 7 pines:
\begin{itemize}[leftmargin=2em]
  \item 3.3V
  \item GND
  \item MOSI (GPIO10)
  \item SCLK (GPIO11)
  \item CE (GPIO8)
  \item DC (GPIO18)
  \item RST (GPIO24)
\end{itemize}

% =====================================================================
\chapter{Layout del PCB}

\section{Tamaño y forma}

\begin{itemize}[leftmargin=2em]
  \item Tamaño: $\sim 90 \times 90$ mm (compatible con estandar PC-104).
  \item Forma: cuadrado con 4 esquinas redondeadas.
  \item 4 agujeros de montaje M3 en las esquinas.
  \item Conectores PC-104 en los bordes (los 104 pines estandar).
\end{itemize}

\section{Stackup recomendado}

\textbf{4 layers}:
\begin{enumerate}[leftmargin=2em]
  \item \textbf{Layer 1 (top)}: signal + componentes.
  \item \textbf{Layer 2}: ground plane continuo.
  \item \textbf{Layer 3}: power plane (dividido por rails 5V, 3.3V, etc.).
  \item \textbf{Layer 4 (bottom)}: signal de retorno + componentes pasivos.
\end{enumerate}

\section{Reglas de routing}

\begin{itemize}[leftmargin=2em]
  \item \textbf{Track width minimum}: 0.15 mm (6 mil) para signal.
  \item \textbf{Track width power}: 0.5 mm para 5V/3.3V; 1.0 mm para
        bus principal.
  \item \textbf{Clearance}: 0.2 mm minimo.
  \item \textbf{Vias}: 0.4 mm hole, 0.8 mm pad.
  \item \textbf{Impedancia controlada} en CSI-2 (90 $\Omega$ diferencial).
  \item \textbf{Sin antenas}: evitar tracks largos paralelos.
\end{itemize}

\section{Componentes criticos}

\textbf{Cerca del CM5}:
\begin{itemize}[leftmargin=2em]
  \item Capacitores de bulk 10 uF en cada rail.
  \item 100 nF en cada pin de power del CM5.
\end{itemize}

\textbf{Cerca del INA219}:
\begin{itemize}[leftmargin=2em]
  \item Shunt 0.1 $\Omega$ con SMD 0805 de 1\% de tolerancia.
  \item Bypass 100 nF en VCC.
\end{itemize}

\textbf{Cerca de los conectores externos}:
\begin{itemize}[leftmargin=2em]
  \item TVS (Transient Voltage Suppressor) en cada pin externo para ESD.
\end{itemize}

% =====================================================================
\chapter{BOM (Bill of Materials)}

Lista preliminar de componentes:

\begin{center}
\small
\begin{tabular}{p{4.5cm}p{4cm}p{3cm}p{1.5cm}}
\toprule
\textbf{Componente} & \textbf{Part number} & \textbf{Proveedor} & \textbf{Costo} \\
\midrule
CM5 (4 GB/32 GB)              & SC1230               & RPi Trading & \$60 \\
DF40 mezzanine (2x)            & Hirose DF40C-100DS-0.4V & Hirose      & \$8 \\
$\mu$Module 5V/6A regulator    & LTM4631              & Analog Devices & \$25 \\
Diodo Schottky 5A              & 1N5824               & Vishay      & \$0.50 \\
Fusible PTC 5A reseteable      & MF-MSMF050           & Bourns      & \$1 \\
INA219 modulo                  & INA219AIDR            & TI          & \$3 \\
TXS0108E level shifter         & TXS0108EPWR           & TI          & \$1 \\
Cap. ceramico 100 nF (50x)     & GRM188R71H104KA93D    & Murata      & \$0.05/u \\
Cap. tantalio 10 uF (10x)      & T491A106K016AT        & Kemet       & \$0.30/u \\
Conector FFC 15-pin            & 5025810-09           & Molex       & \$1 \\
Conector header 7-pin OLED      & estandar 2.54 mm     & ---         & \$0.50 \\
Boton SMD tactile              & 1825910-6             & TE Connectivity & \$0.50 \\
LED bicolor SMD                & APBA3025LSGSGCG       & Kingbright  & \$0.50 \\
Resistor 5.6 ohm 1\%           & varios               & ---         & \$0.05 \\
PCB 4-layer (10 unidades)      & JLCPCB / OshPark      & ---         & \$50 \\
\midrule
\textbf{Total estimado}        & --- & --- & \textbf{\$155} \\
\bottomrule
\end{tabular}
\end{center}

\section{Donde comprar}

\begin{itemize}[leftmargin=2em]
  \item \textbf{Mouser}: el lider para componentes electronicos. Envia a Peru.
  \item \textbf{Digikey}: similar, a veces mejor stock.
  \item \textbf{LCSC}: barato, China, envio lento.
  \item \textbf{Arrow}: para vol\'umenes grandes.
  \item \textbf{JLCPCB}: PCB con assembly opcional, calidad aceptable.
  \item \textbf{OshPark}: PCB premium, mas caro.
\end{itemize}

% =====================================================================
\chapter{Bring-up del carrier}

Procedimiento despues de recibir el PCB ensamblado:

\section{Step 1: inspeccion visual}

\begin{itemize}[leftmargin=2em]
  \item Microscopio para soldaduras: solder bridges, joints frios.
  \item Verificar componentes en sus footprints correctos.
  \item Verificar polaridad de capacitores polarizados.
  \item Verificar orientacion de diodos.
\end{itemize}

\section{Step 2: tests sin CM5 montado}

Antes de poner el CM5:

\begin{itemize}[leftmargin=2em]
  \item Verificar continuidad de GND con multimetro.
  \item Verificar que NO hay corto entre GND y rails (5V, 3V3).
  \item Alimentar con fuente externa a 5V, $\sim$100 mA limite.
  \item Verificar que el regulador entrega 5V estable.
  \item Verificar que INA219 responde via I2C.
\end{itemize}

\section{Step 3: montar el CM5}

\begin{itemize}[leftmargin=2em]
  \item Alinear el CM5 con los conectores DF40.
  \item Presionar firme y uniformemente.
  \item Verificar que esta a 1.5 mm de stack height (no se quedo
        atorado).
  \item Conectar el cable Hostboot via micro-USB / pogo pin.
\end{itemize}

\section{Step 4: flashear la eMMC}

Para el primer boot, hay que flashear el OS en la eMMC:

\begin{lstlisting}
# En una PC linux con el CM5 conectado por USB
# Boton de boot mantenido para entrar a USB device mode
sudo rpiboot

# La eMMC aparece como /dev/sda
# Usar Raspberry Pi Imager para flashear
\end{lstlisting}

\section{Step 5: primer boot}

\begin{itemize}[leftmargin=2em]
  \item Verificar LED de status.
  \item Conectar HDMI debug (si carrier tiene salida) o serial console.
  \item Verificar boot sequence en logs.
  \item Verificar que detecta camara y OLED.
\end{itemize}

\section{Step 6: instalar el daemon cubesat/}

\begin{lstlisting}
cd /opt
git clone https://github.com/JaminYC/CubeSat-EdgeAI-Payload.git
cd CubeSat-EdgeAI-Payload
sudo bash cubesat/install.sh
\end{lstlisting}

% =====================================================================
\chapter{Validacion ambiental}

Una vez bring-up exitoso, hay que validar contra los ambientes de
mision.

\section{Test electrico}

\begin{itemize}[leftmargin=2em]
  \item Variar voltaje de entrada 4.5V a 5.5V, verificar operacion estable.
  \item Test de cortocircuito en salidas (sin daño persistente).
  \item Test de ESD en conectores externos.
  \item Test de inrush current con osciloscopio.
\end{itemize}

\section{Test funcional}

\begin{itemize}[leftmargin=2em]
  \item 100 scans seguidos sin reboot.
  \item I2C OBC bidireccional, 1000 comandos.
  \item Downlink chunked, validar checksums.
  \item Test de SAFE\_MODE: entrada y salida.
  \item Test de OTA: rollback en caso de fallo.
\end{itemize}

\section{Test ambiental}

Como detallado en el manual M7 (Testing):
\begin{itemize}[leftmargin=2em]
  \item TVAC: ciclos -20 a +60 °C, 3 ciclos minimo.
  \item Vibracion: 3 g RMS en cada eje, 1 minuto.
  \item Shock: 1500 g, 0.5 ms (simulando despliegue).
  \item Vacio: $10^{-5}$ Torr durante 24 h (outgassing).
\end{itemize}

\begin{cuidado}
Despues de cada test ambiental, repetir el test funcional para verificar
que no hubo degradacion latente.
\end{cuidado}

% =====================================================================
\chapter{Diferencias con el Pi 5 de desarrollo}

Cosas que tu codigo actual deberia adaptar:

\section{eMMC en lugar de SD}

\begin{itemize}[leftmargin=2em]
  \item No hay slot fisico. La eMMC es soldada.
  \item Para flashear: rpiboot via USB.
  \item Para acceder: como cualquier disco interno.
  \item Backup: copiar imagen completa con \texttt{dd}.
\end{itemize}

\section{No hay Ethernet ni WiFi}

\begin{itemize}[leftmargin=2em]
  \item No se puede SSH durante el vuelo (obvio).
  \item Para debug en lab: el carrier puede tener conector RJ-45 si lo
        decidiste; sino, solo serial console.
\end{itemize}

\section{Distinto pinout del header GPIO}

\begin{itemize}[leftmargin=2em]
  \item Los 40 pines del GPIO siguen iguales en numero BCM.
  \item Pero fisicamente van por el mezzanine, no por header.
\end{itemize}

\section{Sin USB physical}

\begin{itemize}[leftmargin=2em]
  \item Los 4 USB siguen disponibles via mezzanine.
  \item Vos decides si exponerlos en el carrier o no.
\end{itemize}

% =====================================================================
\chapter{Software igual, hardware diferente}

Lo bueno: el sistema operativo (Raspberry Pi OS Bookworm) es \textbf{identico}
en Pi 5 y CM5. Tu pipeline IA corre sin cambios.

\section{Cambios menores en el codigo}

\begin{itemize}[leftmargin=2em]
  \item \texttt{cubesat/paths.py}: cambiar rutas si decides estructura
        diferente (probablemente no cambia).
  \item Verificar que \texttt{pigpio} sigue funcionando (si).
  \item \texttt{cubesat/install.sh}: agregar paso de detectar eMMC vs SD.
  \item Documentacion de boot: actualizar instrucciones.
\end{itemize}

\section{Cambios en configuracion}

\begin{itemize}[leftmargin=2em]
  \item \texttt{/boot/firmware/config.txt}: similar, agregar
        \texttt{dtparam=eeprom\_write\_protect=on} para evitar escritura
        accidental.
  \item Desactivar HDMI: \texttt{hdmi\_force\_hotplug=0}.
  \item Desactivar WiFi/BT (si la version trae): \texttt{dtoverlay=disable-wifi},
        \texttt{dtoverlay=disable-bt}.
\end{itemize}

% =====================================================================
\chapter{Cronograma de migracion}

\begin{itemize}[leftmargin=2em]
  \item \textbf{Semanas 1-2}: comprar CM5 + componentes. Diseñar
        esquematico en KiCad.
  \item \textbf{Semanas 3-4}: layout del PCB. Revision interna.
  \item \textbf{Semana 5}: enviar a fabricar (JLCPCB con assembly).
  \item \textbf{Semanas 6-7}: esperar PCBs. Mientras tanto, preparar
        test plan.
  \item \textbf{Semana 8}: recibir PCBs. Bring-up.
  \item \textbf{Semanas 9-10}: validacion funcional.
  \item \textbf{Semanas 11-14}: tests ambientales (vibracion, TVAC).
  \item \textbf{Semana 15}: integracion al INTISAT.
\end{itemize}

\section{Costo total estimado}

\begin{center}
\begin{tabular}{ll}
\toprule
\textbf{Item} & \textbf{Costo (USD)} \\
\midrule
Componentes para 1 unidad      & 155 \\
PCB (10 unidades)              & 50 \\
Assembly por JLCPCB            & 100 \\
Envio internacional            & 100 \\
Herramientas (estacion de soldadura, microscopio) & 300 (una vez) \\
Aranceles (Peru, 18\%)          & 60 \\
\midrule
\textbf{Total primera iteracion} & \textbf{$\sim$ 765 USD} \\
\bottomrule
\end{tabular}
\end{center}

\section{Numero de iteraciones}

Tipicamente:
\begin{itemize}[leftmargin=2em]
  \item Iteracion 1 (prototipo): tiene bugs menores. Costo $\sim$ \$765.
  \item Iteracion 2 (revisado): casi listo. Costo $\sim$ \$500 (PCBs solo).
  \item Iteracion 3 (qual model): pasa todos los tests. Costo $\sim$ \$500.
  \item Iteracion 4 (flight model): el que vuela. Costo $\sim$ \$500.
\end{itemize}

\textbf{Total proyectado}: \$2200 USD para llegar al flight model.

% =====================================================================
\appendix
\chapter{Apendice --- Plantilla KiCad}

Estructura recomendada del proyecto:

\begin{lstlisting}
intisat-payload-carrier/
|-- intisat-payload-carrier.kicad_pro
|-- intisat-payload-carrier.kicad_sch
|-- intisat-payload-carrier.kicad_pcb
|-- schematics/
|   |-- power.kicad_sch
|   |-- cm5.kicad_sch
|   |-- io_obc.kicad_sch
|   |-- io_camera.kicad_sch
|   |-- io_oled.kicad_sch
|   `-- debug.kicad_sch
|-- libraries/
|   |-- intisat-symbols.kicad_sym
|   `-- intisat-footprints.pretty/
|-- bom/
|   `-- BOM-rev1.csv
|-- outputs/
|   |-- gerbers/
|   |-- 3d-step/
|   `-- docs/
`-- README.md
\end{lstlisting}

\chapter{Apendice --- Pinout DF40 del CM5 (resumen)}

Para el pinout completo ver datasheet oficial. Pines mas usados:

\begin{center}
\small
\begin{tabular}{lll}
\toprule
\textbf{Pin DF40} & \textbf{Funcion} & \textbf{Notas} \\
\midrule
1, 2         & GND                                 & soldar a ground plane \\
3, 5         & 5V\_IN                              & alimentacion principal \\
7, 9         & 3V3\_IN (no usar como output)       & 3.3V para perifericos \\
J1-100       & misc GPIO + buses dedicados         & ver tabla completa \\
J2-1, J2-2   & GND                                 & --- \\
J2-3-100     & otros buses + control               & --- \\
\bottomrule
\end{tabular}
\end{center}

\chapter{Apendice --- Diferencias entre versiones del CM5}

\begin{center}
\begin{tabular}{lll}
\toprule
\textbf{Version} & \textbf{Cambio} & \textbf{Fecha} \\
\midrule
CM5             & version original   & 2024 \\
CM5 Lite        & sin eMMC           & 2024 \\
CM5 + WiFi      & con CYW43455       & 2024 \\
\bottomrule
\end{tabular}
\end{center}

\backmatter

\chapter*{Bibliografia}
\addcontentsline{toc}{chapter}{Bibliografia}

\begin{enumerate}[leftmargin=2em]
\item Raspberry Pi Foundation. \emph{Compute Module 5 Datasheet}. \\
  \url{https://datasheets.raspberrypi.com/cm5/cm5-datasheet.pdf}, 2024.
\item Raspberry Pi Foundation. \emph{CM5 IO Board Datasheet}. \\
  \url{https://datasheets.raspberrypi.com/cm5/cm5-ioboard-datasheet.pdf}, 2024.
\item Analog Devices. \emph{LTM4631 Datasheet}. \\
  \url{https://www.analog.com/en/products/ltm4631.html}, 2020.
\item Texas Instruments. \emph{TXS0108E Datasheet}. \\
  \url{https://www.ti.com/product/TXS0108E}.
\item Texas Instruments. \emph{INA219 Datasheet}. \\
  \url{https://www.ti.com/product/INA219}.
\item Hirose Electric. \emph{DF40 Series Datasheet}. \\
  \url{https://www.hirose.com/en/product/series/DF40}.
\item KiCad. \emph{Official Documentation}. \\
  \url{https://docs.kicad.org/}.
\end{enumerate}

\end{document}
